\documentclass[11pt]{article}
\usepackage{fontspec}
\setmainfont[
  Path=fonts/,
  Extension=.ttf,
  UprightFont=*-400,
  BoldFont=*-700
]{NotoSerifSC}
\setsansfont[
  Path=fonts/,
  Extension=.ttf,
  UprightFont=*-400,
  BoldFont=*-700
]{NotoSerifSC}
\XeTeXlinebreaklocale "zh"
\XeTeXlinebreakskip=0pt plus 1pt
\usepackage[a4paper,margin=1.70cm]{geometry}
\usepackage{amsmath,amssymb,mathtools,bm}
\usepackage{booktabs,longtable,array}
\usepackage{enumitem}
\setlength{\parindent}{2em}
\setlength{\parskip}{0.35em}
\setlist{nosep,leftmargin=2em}
\allowdisplaybreaks[3]
\numberwithin{equation}{section}
\pagestyle{plain}

\newcommand{\R}{\mathbb{R}}
\newcommand{\E}{\mathbb{E}}
\newcommand{\Pp}{\mathbb{P}}
\newcommand{\1}{\mathbf{1}}
\newcommand{\T}{\mathsf{T}}
\newcommand{\F}{\mathrm{F}}
\newcommand{\cN}{\mathcal{N}}
\newcommand{\cL}{\mathcal{L}}
\newcommand{\cS}{\mathcal{S}}
\newcommand{\cH}{\mathcal{H}}
\newcommand{\cA}{\mathcal{A}}
\newcommand{\cM}{\mathcal{M}}
\newcommand{\cT}{\mathcal{T}}
\newcommand{\diag}{\operatorname{diag}}
\newcommand{\tr}{\operatorname{tr}}
\newcommand{\Var}{\operatorname{Var}}
\newcommand{\Cov}{\operatorname{Cov}}
\newcommand{\softmax}{\operatorname{softmax}}
\newcommand{\KL}{D_{\mathrm{KL}}}
\newcommand{\sg}{\operatorname{sg}}
\begin{document}
    \begin{center}
      {\LARGE\bfseries 1-2106.09685v2-Basis-LoRA}
    \end{center}
    \vspace{0.8em}

    \section{符号与维度}
    \begin{longtable}{>{$}l<{$} >{$}l<{$} p{10cm}}
    \toprule
    \text{符号} & \text{维度} & \text{含义}\\
    \midrule
    W_0 & \R^{m\times n} & 冻结预训练权重。\\
A,B & \R^{r\times n},\R^{m\times r} & 低秩下投影与上投影。\\
r & \mathbb N & 适配器秩。\\
x,\delta & \R^n,\R^m & 输入和输出端反向梯度。\\
c & \R & 低秩更新缩放系数。\\
    \bottomrule
    \end{longtable}

    所有向量默认是列向量；未特别说明时，期望同时对数据、噪声和采样时刻取值。下文只展开论文中承担方法逻辑的公式，并把所需假设写在推导发生的位置。

\section{低秩更新的完整矩阵微分}

令冻结权重 $W_0\in\R^{m\times n}$，低秩更新
\begin{equation}
 W=W_0+cBA,\qquad
 B\in\R^{m\times r},\quad A\in\R^{r\times n},\quad c=\alpha/r.
\end{equation}
对一个输入 $x\in\R^n$，
\begin{equation}
 h=W_0x+cB(Ax).
\end{equation}
令 $y=Ax\in\R^r$、输出梯度 $g=\partial\cL/\partial h\in\R^m$。
微分为
\begin{equation}
 dh=c\,dB\,y+cB\,dA\,x+cBA\,dx+W_0dx.
\end{equation}
由 $d\cL=g^Tdh$ 以及
$a^TdMb=\langle ab^T,dM\rangle_\F$，逐项得到
\begin{align}
 \boxed{\nabla_B\cL=c,g,y^T=c,g,x^TA^T,}\\
 \boxed{\nabla_A\cL=c,B^Tg,x^T,}\\
 \boxed{\nabla_x\cL=W_0^Tg+cA^TB^Tg.}
\end{align}
若批输入按行堆成 $X\in\R^{b\times n}$，上游梯度为
$G\in\R^{b\times m}$，则
\begin{align}
 \nabla_B\cL&=cG^TXA^T,\\
 \nabla_A\cL&=cB^TG^TX.
\end{align}

\subsection{rank-1 专家分解}

写 $B=[b_1,\ldots,b_r]$、$A$ 的第 $i$ 行为 $a_i^T$，则
\begin{equation}
 BA=\sum_{i=1}^{r}b_i a_i^T,
 \qquad BAx=\sum_{i=1}^{r}b_i(a_i^Tx).
\end{equation}
每一项均为秩至多 $1$ 的更新。对单样本，
\begin{equation}
 \nabla_{b_i}\cL=c(a_i^Tx)g,\qquad
 \nabla_{a_i}\cL=c(b_i^Tg)x.
\end{equation}
于是列 $b_i$ 的更新强度由标量路由量 $a_i^Tx$ 控制；这也是把 rank 分量解释为隐式专家的数学依据。

\section{初始化时两因子的非对称动力学}

常见初始化为 $B_0=0$、$A_0$ 随机。此时
\begin{equation}
 \Delta W_0=cB_0A_0=0,
\end{equation}
所以前向不改变基座。梯度却满足
\begin{equation}
 \nabla_B\cL\big|_{0}=cG^TXA_0^T,\qquad
 \nabla_A\cL\big|_{0}=cB_0^TG^TX=0.
\end{equation}
第一步梯度下降
\begin{equation}
 B_1=-\eta cG^TXA_0^T,\qquad A_1=A_0,
\end{equation}
因此第二步开始 $A$ 才获得非零梯度。若反过来令 $A_0=0$、$B_0$ 随机，
则第一步只有 $A$ 更新。两种初始化前向等价，但早期优化动力学不同。

还有尺度不识别性：任意 $s\ne0$ 满足
\begin{equation}
 BA=(sB)(s^{-1}A).
\end{equation}
然而带独立权重衰减时
\begin{equation}
 \|sB\|_F^2+\|s^{-1}A\|_F^2
 =s^2\|B\|_F^2+s^{-2}\|A\|_F^2,
\end{equation}
最优平衡尺度由求导得到
\begin{equation}
 s^4=\frac{\|A\|_F^2}{\|B\|_F^2}.
\end{equation}
故因子正则化会隐式选择分解，即使乘积 $BA$ 相同。

\section{最佳秩约束近似}

设目标全量更新的奇异值分解为
\begin{equation}
 \Delta W^*=U\Sigma V^T
 =\sum_{i=1}^{q}\sigma_i u_iv_i^T,
 \qquad\sigma_1\ge\cdots\ge\sigma_q\ge0.
\end{equation}
截断矩阵
\begin{equation}
 \Delta W_r=\sum_{i=1}^{r}\sigma_i u_iv_i^T
\end{equation}
满足 Eckart--Young--Mirsky 结论
\begin{equation}
 \min_{\operatorname{rank}(M)\le r}
 \|\Delta W^*-M\|_F^2
 =\sum_{i>r}\sigma_i^2.
\end{equation}
一种直接证明如下。因 Frobenius 范数在正交变换下不变，令
$C=U^TMV$，则 $\operatorname{rank}(C)\le r$ 且
\begin{equation}
 \|\Delta W^*-M\|_F^2=\|\Sigma-C\|_F^2.
\end{equation}
任意秩不超过 $r$ 的 $C$ 至多精确保留 $r$ 个相互正交奇异方向；
由奇异值的变分刻画或 Pythagoras 分解，剩余平方能量至少为
$\sum_{i>r}\sigma_i^2$。取
$C=\diag(\sigma_1,\ldots,\sigma_r,0,\ldots)$ 达到该下界。

谱范数版本为
\begin{equation}
 \min_{\operatorname{rank}(M)\le r}\|\Delta W^*-M\|_2=\sigma_{r+1}.
\end{equation}
这只是表达能力上界；训练能否找到截断 SVD 解还取决于非凸因子化、数据和优化器。

\section{子空间限制与梯度投影}

若 $A$ 冻结，只训练 $B$，可达更新集合为
\begin{equation}
 \cS_A=\{BA:B\in\R^{m\times r}\}.
\end{equation}
其中每一行都位于 $A$ 的行空间。若 $A$ 满行秩，其行空间正交投影为
\begin{equation}
 P_A=A^T(AA^T)^{-1}A.
\end{equation}
给定任意目标更新 $M$，最小二乘问题
\begin{equation}
 \min_B\|M-BA\|_F^2
\end{equation}
的正规方程为
\begin{equation}
 (BA-M)A^T=0,\qquad
 B^*=MA^T(AA^T)^{-1}.
\end{equation}
故最佳可达更新
\begin{equation}
 B^*A=MP_A,\qquad
 \min_B\|M-BA\|_F^2=\|M(I-P_A)\|_F^2.
\end{equation}
这把“冻结随机下投影”的表达损失精确写成目标在随机行空间外的能量。

\section{随机投影的尺度与集中}

设 $A_{ij}$ 独立、均值零、方差 $1/r$。对固定 $x\in\R^n$，
\begin{align}
 \E\|Ax\|_2^2
 &=\sum_{i=1}^{r}\E\left(\sum_jA_{ij}x_j\right)^2\\
 &=\sum_i\sum_j\E[A_{ij}^2]x_j^2
 =\|x\|_2^2.
\end{align}
若 $A_{ij}\sim\cN(0,1/r)$，则
\begin{equation}
 \frac{\|Ax\|_2^2}{\|x\|_2^2}
 \sim\frac1r\chi_r^2.
\end{equation}
利用卡方尾界，对 $0<\varepsilon<1$，
\begin{equation}
 \Pp\left(
 \left|\|Ax\|_2^2-\|x\|_2^2\right|
 \ge\varepsilon\|x\|_2^2\right)
 \le2\exp\left(-\frac r8\varepsilon^2\right).
\end{equation}
对有限集合 $N$ 个向量做并合界，若
\begin{equation}
 r\ge\frac8{\varepsilon^2}\log\frac{2N}{\delta},
\end{equation}
则所有向量同时保持范数的概率至少为 $1-\delta$。对点对差向量应用同一结论即得到距离保持。

稀疏随机投影若每个元素以概率 $\rho$ 非零，则为保持同样二阶矩，非零幅值应按
$1/\sqrt{\rho r}$ 缩放：
\begin{equation}
 A_{ij}=\begin{cases}
 +1/\sqrt{\rho r},&\rho/2,\\
 -1/\sqrt{\rho r},&\rho/2,\\
 0,&1-\rho.
 \end{cases}
\end{equation}
此时 $\E[A_{ij}^2]=1/r$，但四阶矩为
$\E[A_{ij}^4]=1/(\rho r^2)$；$\rho$ 越小，尾部越重，获得同样集中精度通常需要更大 $r$。

\section{任务更新的干扰量}

两个任务的更新 $U_i=B_iA_i$、$U_j=B_jA_j$ 的 Frobenius 内积为
\begin{align}
 \langle U_i,U_j\rangle_F
 &=\tr((B_iA_i)^TB_jA_j)\\
 &=\tr(B_i^TB_jA_jA_i^T).
\end{align}
因此
\begin{equation}
 |\langle U_i,U_j\rangle_F|
 \le\|B_i^TB_j\|_F\|A_jA_i^T\|_F.
\end{equation}
随机下投影近似正交只控制第二个因子；若 $B_i,B_j$ 范数无界，绝对干扰仍可任意大。
对加权合并 $U=\sum_iw_iU_i$，
\begin{equation}
 \|U\|_F^2
 =\sum_iw_i^2\|U_i\|_F^2
 +2\sum_{i<j}w_iw_j\langle U_i,U_j\rangle_F.
\end{equation}
交叉项就是任务干扰对合并范数的精确贡献，而不是抽象的“方向冲突”。

\section{参数量、训练内存与合并}

全量矩阵参数量为 $mn$，低秩因子参数量为
\begin{equation}
 r(m+n),
\end{equation}
压缩要求 $r<mn/(m+n)$。训练时只有可训练因子需要梯度和优化器状态；
冻结基座仍需存储权重并参与前向、反向的输入梯度计算。

推理前可合并
\begin{equation}
 W_{\rm merged}=W_0+cBA.
\end{equation}
合并后对单个适配器不增加矩阵乘次数，但丢失了动态切换不同适配器的便利。
若基座以低比特存储，通常必须先在计算精度中反量化再加 $BA$；直接在整数码空间相加并不等价。

\section{光滑损失的下降引理}

若 $L$ 的梯度是 $\beta$-Lipschitz：
\begin{equation}
 \|\nabla L(x)-\nabla L(y)\|
 \le\beta\|x-y\|,
\end{equation}
则沿线段积分
\begin{align}
 L(y)-L(x)
 &=\int_0^1\nabla L(x+t(y-x))^T(y-x)dt\\
 &=\nabla L(x)^T(y-x)+R,
\end{align}
余项满足
\begin{align}
 |R|
 &\le\int_0^1\beta t\|y-x\|^2dt
 =\frac\beta2\|y-x\|^2.
\end{align}
故下降引理
\begin{equation}
 L(y)\le L(x)+\nabla L(x)^T(y-x)
 +\frac\beta2\|y-x\|^2.
\end{equation}
取梯度步 $y=x-\eta\nabla L(x)$：
\begin{equation}
 L(y)\le L(x)-\eta\left(1-\frac{\beta\eta}{2}\right)
 \|\nabla L(x)\|^2.
\end{equation}
当 $0<\eta<2/\beta$ 时每个精确梯度步下降；深网中 $\beta$ 未知且随位置变化，
这仍解释了学习率过大造成不稳定的局部机制。

\section{二阶近似与曲率方向}

在 $x$ 附近 Taylor 展开
\begin{equation}
 L(x+\Delta)=L(x)+g^T\Delta
 +\frac12\Delta^TH\Delta+O(\|\Delta\|^3).
\end{equation}
若 $H=U\diag(\lambda_i)U^T$，在本征方向 $u_i$ 上步长 $a$ 的变化
\begin{equation}
 \Delta L\approx a g_i+\frac12\lambda_i a^2.
\end{equation}
$\lambda_i>0$ 为局部碗形，$\lambda_i<0$ 是负曲率逃离鞍点方向。
加阻尼的 Newton 步
\begin{equation}
 \Delta=-(H+\lambda I)^{-1}g
\end{equation}
在特征方向缩放为 $-g_i/(\lambda_i+\lambda)$，避免接近零或负特征值造成巨大步长。

\section{链式 Jacobian 与条件数}

复合映射 $y=f_L\circ\cdots\circ f_1(x)$ 的 Jacobian
\begin{equation}
 J_{y\leftarrow x}=J_LJ_{L-1}\cdots J_1.
\end{equation}
谱范数上界
\begin{equation}
 \|J_{y\leftarrow x}\|_2\le\prod_{\ell=1}^{L}\|J_\ell\|_2.
\end{equation}
若各层最小奇异值存在，
\begin{equation}
 \sigma_{\min}(J_{y\leftarrow x})
 \ge\prod_\ell\sigma_{\min}(J_\ell).
\end{equation}
因此条件数至多乘法累积：
\begin{equation}
 \kappa(J_{y\leftarrow x})
 \le\prod_\ell\kappa(J_\ell).
\end{equation}
低维映射、递归模块和物理传播都面临同一问题：表示存在不代表反向梯度数值良态。

\section{随机梯度的偏差与方差}

设样本梯度 $g(\theta;\xi)$ 满足
\begin{equation}
 \E_\xi g(\theta;\xi)=\nabla L(\theta),
 \qquad \E\|g-\nabla L\|^2\le\sigma^2.
\end{equation}
独立批量 $B$ 的平均
\begin{equation}
 \bar g_B=\frac1B\sum_{i=1}^{B}g_i
\end{equation}
满足
\begin{equation}
 \E\bar g_B=\nabla L,\qquad
 \E\|\bar g_B-\nabla L\|^2\le\frac{\sigma^2}{B}.
\end{equation}
若样本相关，协方差项出现：
\begin{equation}
 \Var(\bar g_B)=\frac1{B^2}\left[
 \sum_i\Var(g_i)+\sum_{i\ne j}\Cov(g_i,g_j)
 \right].
\end{equation}
扩大批量只有在交叉协方差不主导时才近似按 $1/B$ 降噪。

若使用有偏估计 $\E\widehat g=\nabla L+b$，均方误差分解
\begin{equation}
 \E\|\widehat g-\nabla L\|^2
 =\|b\|^2+\tr\Cov(\widehat g).
\end{equation}
直通估计、停止梯度目标和截断反传通常以可控偏差换低方差或低内存。

\section{Adam 的尺度归一化}

逐坐标 Adam 更新
\begin{align}
 m_t&=\beta_1m_{t-1}+(1-\beta_1)g_t,\\
 v_t&=\beta_2v_{t-1}+(1-\beta_2)g_t^2,\\
 \widehat m_t&=m_t/(1-\beta_1^t),\\
 \widehat v_t&=v_t/(1-\beta_2^t),\\
 \theta_{t+1}&=\theta_t-\eta
 \frac{\widehat m_t}{\sqrt{\widehat v_t}+\epsilon}.
\end{align}
若梯度长期乘常数 $c>0$，忽略 $\epsilon$ 时
$\widehat m$ 乘 $c$、$\sqrt{\widehat v}$ 也乘 $c$，更新近似不变。
但瞬态、符号变化和 $\epsilon$ 会破坏精确尺度不变性。

AdamW 把权重衰减解耦：
\begin{equation}
 \theta_{t+1}=(1-\eta\lambda)\theta_t
 -\eta\frac{\widehat m_t}{\sqrt{\widehat v_t}+\epsilon}.
\end{equation}
这不同于把 $\lambda\theta$ 加入梯度后再经过自适应分母；后者会按坐标二阶矩缩放正则化。

\section{范数约束与投影梯度}

约束集合 $\mathcal C$ 上的投影梯度
\begin{equation}
 \theta^+=\Pi_{\mathcal C}(\theta-\eta g),
 \qquad
 \Pi_{\mathcal C}(x)=\arg\min_{y\in\mathcal C}\|x-y\|^2.
\end{equation}
对欧氏球 $\mathcal C=\{x:\|x\|\le R\}$，
\begin{equation}
 \Pi_{\mathcal C}(x)=
 \begin{cases}
 x,&\|x\|\le R,\\
 Rx/\|x\|,&\|x\|>R.
 \end{cases}
\end{equation}
对流形约束，局部一阶步先投影到切空间，再用 retraction 回到流形；
仅把参数写成 $g(z)$ 相当于用参数化隐式满足约束，但会引入 $J_g^TJ_g$ 的坐标度量。

\section{Lipschitz 误差的层间传播}

真实模块 $f_\ell$、近似模块 $\widehat f_\ell$ 满足
\begin{equation}
 \|f_\ell(x)-\widehat f_\ell(x)\|\le\epsilon_\ell,
 \qquad
 \|f_\ell(x)-f_\ell(y)\|\le L_\ell\|x-y\|.
\end{equation}
令真实、近似中间状态差 $e_\ell$。三角不等式给出
\begin{align}
 e_{\ell+1}
 &=\|f_\ell(h_\ell)-\widehat f_\ell(\widehat h_\ell)\|\\
 &\le L_\ell e_\ell+\epsilon_\ell.
\end{align}
递归展开
\begin{equation}
 e_L\le
 \left(\prod_{j=0}^{L-1}L_j\right)e_0
 +\sum_{\ell=0}^{L-1}
 \left(\prod_{j=\ell+1}^{L-1}L_j\right)\epsilon_\ell.
\end{equation}
早期层误差乘更多后续 Lipschitz 因子，所以量化、低秩近似或局部训练的相同单层误差，
放在不同深度并不等价。

\section{复杂度必须区分四个量}

参数量、FLOPs、激活内存和通信量是不同指标。以线性层
$Y=XW$、$X\in\R^{B\times n}$、$W\in\R^{n\times m}$ 为例：
\begin{align}
 \text{参数量}&=nm,\\
 \text{前向乘加量}&\asymp Bnm,\\
 \text{主要激活量}&=B(n+m),\\
 \text{数据并行梯度通信}&\asymp nm.
\end{align}
低秩因子可同时减少参数与乘法，但动态稀疏若实现仍做稠密 kernel，参数/FLOPs 公式不会自动转化为墙钟加速。
冻结参数减少梯度和优化器状态，却不一定减少前向计算或输入梯度计算。

\end{document}
